% sections/01-introduction.tex

\section{Introduction}
\label{sec:intro}

Visual complexity and aesthetic perception are among the most fundamental yet elusive qualities in art and design images. For decades, researchers in visual communication, design studies, and computer vision have sought computational proxies that capture how humans perceive and evaluate visual content \citep{berlyne1971aesthetics, leder2004model}. Automatic assessment could support interface design evaluation, artwork curation, content recommendation, and design quality control \citep{zhang2020survey}. However, two persistent obstacles limit progress.

First, large-scale human-rated datasets are prohibitively expensive, especially when multiple perceptual dimensions are required. As a result, many studies develop models on computational pseudo-labels---scores generated by theory-derived formulas or teacher metrics---and then claim human-level performance without an independent audit. This risks circular validation: the model may only reconstruct the pseudo-label formula rather than predict genuine human judgments \citep{damasio2022visual}. Second, deep-learning approaches, while powerful on large annotated datasets such as AVA \citep{murray2012ava}, require substantial resources and offer limited interpretability, creating a barrier for design researchers who need to understand \textit{why} a visual property influences a perceptual score \citep{li2023interpretable}.

This paper addresses the gap between cheap computational proxies and costly human validation by asking how much human evidence is actually necessary to audit a lightweight, interpretable perceptual proxy. Our research questions are:

\begin{itemize}
    \item \textbf{RQ1 -- Proxy construction:} Can compact, design-theory-grounded visual features produce stable and computationally efficient proxy scores across diverse art and design image categories?
    \item \textbf{RQ2 -- Human alignment:} How well do frozen proxy-model predictions align with independent human ratings without using those ratings for feature selection or model tuning?
    \item \textbf{RQ3 -- Label efficiency:} How many rated images and how many raters per image are required to obtain stable model-evaluation conclusions?
    \item \textbf{RQ4 -- Interpretability:} Which visual factors consistently influence proxy predictions and human-aligned predictions, and where do the two explanations diverge?
    \item \textbf{RQ5 -- Boundary conditions:} Which perceptual dimensions and image categories cannot be assessed reliably by the proposed lightweight representation?
\end{itemize}

We frame the task as \textit{proxy assessment} rather than direct human-perception prediction. The framework uses 30 core handcrafted visual design features and Ridge regression. All feature definitions, preprocessing choices, and hyperparameters are frozen before the human validation set is inspected. The large-scale dataset (17,337 images) supports feature development, model selection, and robustness analysis; the 100-image, 30-rater blind-evaluation dataset supports reliability analysis, frozen external validation, and carefully separated optional calibration.

The contributions of this work are fourfold:

\begin{enumerate}
    \item \textbf{C1 -- Methodological separation.} We implement an explicitly layered pipeline (large-scale proxy development, frozen human validation, human-efficiency analysis, optional calibration) that prevents the human validation set from influencing model selection.
    \item \textbf{C2 -- Empirical human alignment.} We report frozen external validation results with bootstrap confidence intervals and false-discovery-rate correction. The lightweight proxy correlates with human mean ratings for Complexity ($\rho = 0.773$), Order ($\rho = 0.724$), Beauty ($\rho = 0.565$), and Emotion ($\rho = 0.408$).
    \item \textbf{C3 -- Label-efficiency evidence.} Rater- and image-saturation experiments quantify the human-evaluation budget. In this task, 10--20 raters and 80--100 stratified images yield stable validation conclusions.
    \item \textbf{C4 -- Interpretability and efficiency.} The model is fully linear, so feature importance is given by standardized Ridge coefficients. Inference takes approximately 0.07\,ms per image and the saved model occupies 7.7\,KB.
\end{enumerate}

We also report clear negative findings. Human ratings of Hierarchy are unreliable (ICC(2,1) = 0.025) and the proxy does not claim to capture this dimension. The framework is therefore positioned as a human-light screening tool, not a replacement for human judgment.

The remainder of the paper is organized as follows. Section~\ref{sec:related} reviews related work. Section~\ref{sec:method} describes datasets, proxy targets, features, model selection, and the human-rating protocol. Section~\ref{sec:experiments} details the evaluation protocol. Section~\ref{sec:results} reports the redesigned experiments. Section~\ref{sec:discussion} discusses implications and limitations. Section~\ref{sec:conclusion} concludes.
